\begin{abstract}
本文研究半径为 $1800\ \mathrm m$ 的圆形区域内无线电干扰源的自动定位与清除问题。示向误差、接收距离和发射方向都存在不确定性，我们建立了从\textbf{几何交会}、\textbf{选点规划}到\textbf{多源搜索清除}的递进式模型，并用数值算例和模拟器日志检验其正确性和可执行性。

\textbf{针对问题一，}将每次示向误差区域表示为两个闭半平面的交，证明有效交会所得定位区域为\textbf{凸多边形}；利用\textbf{旋转卡壳法}求取最远顶点对作为区域直径，并以全部顶点是否位于直径圆内作为\textbf{覆盖判据}，从而给出可直接实现的\textbf{交会定位算法}。

\textbf{针对问题二，}用\textbf{外切正多边形}构造源位置可行域的保守外近似，用\textbf{内接正多边形}构造\textbf{保证接收域}，并以\textbf{最小包围圆半径}衡量定位不确定性。分别采用有限离散综合评分和 \textbf{连续 Fisher 信息搜索}产生候选点，再按\textbf{最坏场景定位半径}与行动时间统一重新评分。在 $200$ 个配对场景中，连续方案将最大定位不确定性半径均值由 $119.037\ \mathrm m$ 降至 $70.082\ \mathrm m$，但综合评分仅在 $76$ 个场景中更优，说明\textbf{选点}需同时权衡定位收益与移动代价。

\textbf{针对问题三，}在半径 $960\ \mathrm m$ 的圆周上设置八个\textbf{驻留点}实现全向源的\textbf{保证发现}，并结合定位区域递推、有限次\textbf{补充测向}和开放式\textbf{旅行商动态规划}，形成“\textbf{全域扫描}---\textbf{批量定位}---\textbf{路径清除}”的闭环策略。汇总 $50$ 条模拟日志，共 $648$ 个实际干扰源全部清除，平均定位清除时间为 \textbf{332.496 s/个}。三个正式测试案例全部清除，平均定位清除时间分别为 \textbf{293.036 s}、\textbf{427.066 s} 和 \textbf{289.209 s}。

\textbf{针对问题四，}针对定向源的半圆可见性，构造边长 $985\ \mathrm m$ 的\textbf{平移等边三角网}，以\textbf{多方向驻留点}保证发现；进一步建立包含源位置、接收半径、发射方向和源类型的\textbf{有限状态集}，利用\textbf{多方向探测组}与\textbf{条件尾部均值}进行\textbf{滚动选点}。汇总 $50$ 条模拟日志，共 $626$ 个实际干扰源全部清除，平均定位清除时间为 \textbf{798.755 s/个}。三个正式测试案例全部清除，平均定位清除时间分别为 \textbf{812.934 s}、\textbf{630.479 s} 和 \textbf{619.302 s}。

检验结果表明，问题一的普通与边界数据共 $1270$ 项检查全部通过，问题三、四在当前测试条件下均达到 $100\%$ 的清除比例。所建模型能够在有界误差和行动反馈条件下完成全向、定向干扰源的\textbf{自动搜索定位与清除}；\textbf{有限离散}和\textbf{分阶段规划}保证了算法可执行，但不构成连续原问题的\textbf{全局最优性}证明。

\keywords{无线电干扰源\quad 交会定位\quad 凸几何\quad 路径规划\quad 滚动决策}
\end{abstract}
